\begin{figure}[htbp]
\centering
\begin{tikzpicture}[fd,
  fd2 ax/.style={lcaxis,gray!60!black},
  fd2 lc/.style={draw=gray!55,line width=0.5pt,dash pattern=on 1.2pt off 1.6pt},
  fd2 warrow/.style={postaction={decorate,decoration={markings,
      mark=at position #1 with {\arrow[xshift=3.2pt]{Stealth[length=6.4pt,width=6pt]}}}}},
  fd2 heavy/.style={draw,solid,line width=1.4pt,fd arrow=#1},
  fd2 heavy/.default=0.5,
  fd2 sheet/.style={fill=gray!13,draw=none},
  fd2 lab/.style={font=\small},
  fd2 note/.style={font=\footnotesize}]
  % geometry: beta = 0.55, t = t' = 1.85; perspective b_perp vector B = (1.7,0.55)
  \foreach \panel/\xs in {a/0, b/7.6} {
  \begin{scope}[xshift=\xs cm]
    \coordinate (O) at (0,0);          \coordinate (OB) at (1.7,0.55);
    \coordinate (P) at (-1.0175,-1.85); \coordinate (PB) at (0.6825,-1.3);
    \coordinate (F) at (-1.0175,1.85);  \coordinate (FB) at (0.6825,2.4);
    % axes and light cone (b_perp = 0)
    \draw[fd2 lc] (-1.75,-1.75) -- (1.75,1.75);  \draw[fd2 lc] (-1.75,1.75) -- (1.75,-1.75);
    \draw[fd2 ax] (-1.75,0) -- (2.75,0) node[below left=1pt and -3pt,font=\footnotesize] {$z$};
    \draw[fd2 ax] (0,-2.25) -- (0,2.95) node[left,font=\footnotesize] {$t$};
    \draw[gray!60!black,line width=0.5pt,dash pattern=on 3pt off 2pt,-{Stealth[length=4pt,width=3.4pt]}] (O) -- (2.45,0.79) node[right,font=\footnotesize] {$\vec b_\perp$};
    \if a\panel
      % two staple-shaped Wilson lines: past (v) and future (v'), closed at -t and t'
      \draw[wilson plain] (OB) -- (PB) -- (P) -- (O) -- (F) -- (FB) -- cycle;
      \path[fd2 warrow=0.55] (P) -- (O);    \path[fd2 warrow=0.6] (O) -- (F);
      \path[fd2 warrow=0.55] (FB) -- (OB);  \path[fd2 warrow=0.6] (OB) -- (PB);
      \path[fd2 warrow=0.55] (F) -- (FB);   \path[fd2 warrow=0.55] (PB) -- (P);
      \node[opvertex] at (O) {};  \node[opvertex] at (OB) {};
      \node[fd2 lab,anchor=east] at (-0.6,-1.0) {$v$};
      \node[fd2 lab,anchor=east] at (-0.6,1.0) {$v'$};
      \node[fd2 note,anchor=north east] at ($(P)+(0.02,-0.04)$) {$-t$};
      \node[fd2 note,anchor=south east] at ($(F)+(0.02,0.04)$) {$t'$};
      \node[fd2 note,anchor=north west] at (0.05,-0.02) {$0$};
      \node[fd2 lab,align=center] at (0.45,-3.1) {$S=\displaystyle\lim_{t,t'\to\infty}\frac{W(t,t',b_\perp,Y,Y')}{\sqrt{Z(2t,b_\perp,Y)\,Z(2t',b_\perp,Y')}}$};
      \node[fd2 note,align=left,anchor=west] at (1.95,1.75) {时序乘积 $\mathcal T$};
      \node[fd2 lab] at (0.45,-4.2) {(a)\ 订书钉形 Wilson 回路};
    \else
      % HQET picture: static colour sources Q (at 0) and Qbar (at b_perp)
      \fill[fd2 sheet] (P) -- (O) -- (F) -- (FB) -- (OB) -- (PB) -- cycle;
      \draw[fd2 heavy=0.55] (P) -- (O);    \draw[fd2 heavy=0.55] (O) -- (F);
      \draw[fd2 heavy=0.5] (FB) -- (OB);   \draw[fd2 heavy=0.55] (OB) -- (PB);
      \draw[gray!60!black,line width=0.6pt,dash pattern=on 2.5pt off 1.5pt] (O) -- (OB);
      \node[opvertex] at (O) {};  \node[opvertex] at (OB) {};
      \node[fd2 lab,anchor=east] at (-0.6,-1.0) {$v$};
      \node[fd2 lab,anchor=east] at (-0.6,1.0) {$v'$};
      \node[fd2 note,anchor=east] at (-1.05,-1.9) {$Q$};
      \node[fd2 note,anchor=west] at (0.72,-1.38) {$\bar Q$};
      \node[fd2 note,anchor=south,inner sep=1pt] at (0.62,0.44) {$J$（$t=0$）};
      \node[fd2 lab,anchor=north] at (-0.15,-2.15) {$|Q\bar Q,\vec b_\perp\rangle_v$};
      \node[fd2 lab,anchor=west] at (0.8,2.6) {${}_{v'}\langle Q\bar Q,\vec b_\perp|$};
      \node[draw=gray!55,fill=gray!4,rounded corners=2pt,font=\footnotesize,align=left,anchor=north west,inner sep=3pt] at (1.75,-0.55)
        {欧氏表示：\\$\ii\tilde v\cdot D_E=D_\tau-\ii\beta D_z$\\（$\tau$：欧氏时间）};
      \node[fd2 lab] at (0.45,-4.2) {(b)\ HQET 形状因子};
    \fi
  \end{scope}}
  \node[fd2 lab,align=center] at (8.05,-3.1) {$S={}_{v'}\langle Q\bar Q,\vec b_\perp|\,J\,|Q\bar Q,\vec b_\perp\rangle_v$};
\end{tikzpicture}
\caption{离光锥 TMD 软函数作为快速色源对的形状因子（$v=\gamma(1,\beta,\mathbf 0_\perp)$、$v'=\gamma'(1,-\beta',\mathbf 0_\perp)$，$\beta=\tanh Y$、$\beta'=\tanh Y'$，坐标 $(t,z,\vec b_\perp)$，$\vec b_\perp$ 方向按透视画出）。
(a)~两个订书钉形 Wilson 回路（$t<0$ 沿 $v$，$t>0$ 沿 $v'$，在 $-t$ 与 $t'$ 处由横向连接闭合）的时序真空期望，除以 $\sqrt{ZZ}$（$Z$ 为矩形回路）后取 $t,t'\to\infty$；箭头为路径方向：$0$ 处沿世界线指向未来（色源 $Q$），$\vec b_\perp$ 处指向过去（$\bar Q$）。
(b)~等价的 HQET 图像：间隔 $b_\perp$ 的色单态静态 $Q\bar Q$ 对以速度 $v$ 运动，在 $t=0$ 被等时的流 $J$ 踢到 $v'$，软函数是 ${}_{v'}\langle Q\bar Q,\vec b_\perp|J|Q\bar Q,\vec b_\perp\rangle_v$，因而有欧氏表示（\cref{thm:yzliu-soft-hqet}，全阶论证 [A]；非微扰的形状因子形式另需\cref{asm:yzliu-heavy-pair}）。
大快度下 $\ln S=(Y+Y')\tilde K+\mathcal D+\order{\ee^{-(Y+Y')}}$，本征软函数 $S_I=\ee^{-\mathcal D}$（VF-YZL-2）。}
\label{fig:yzliu-soft-hqet}
\end{figure}
